\section{制度背景与机制}

\begin{frame}{制度背景：两化融合贯标}
  \begin{itemize}
    \item 两化融合：以信息化推动工业化、以工业化促进信息化。
    \item 管理体系贯标将数字化从零散技术改造转化为可复制、可审计、可持续改进的组织流程。
    \item 截至2025年4月，贯彻两化融合管理体系标准的企业超过7.4万家。
    \item 贯标严格依据国家统一标准，由独立第三方机构现场审查与分级认证。
    \item 地区与年份维度上的分批推进，为交叠双重差分提供政策环境。
  \end{itemize}
\end{frame}

\begin{frame}{政策推进：财政激励与试点示范}
  \begin{figure}
    \centering
    \includegraphics[width=0.82\linewidth,height=0.62\textheight,keepaspectratio]{Subsidy_Trend.pdf}
    \caption{上市公司获得的“两化融合”相关政府补助}
  \end{figure}
\end{frame}

\begin{frame}{宏观事实：固定投资与研发投入分化}
  \begin{figure}
    \centering
    \includegraphics[width=0.82\linewidth,height=0.62\textheight,keepaspectratio]{Macro_RDShare_LeftAxis_Index_RightAxis_1995base_BW.pdf}
    \caption{固定投资与研发投入的相对变化}
  \end{figure}
\end{frame}

\begin{frame}{微观事实：实物投资与研发投入}
  \begin{figure}
    \centering
    \includegraphics[width=0.82\linewidth,height=0.62\textheight,keepaspectratio]{Micro_Investment_Structure_Scissors.pdf}
    \caption{上市企业投资结构的剪刀差}
  \end{figure}
\end{frame}

\begin{frame}{机制：抵押品替代与投资结构}
  \begin{columns}[T]
    \column{0.48\textwidth}
      \textbf{传统约束}
      \begin{itemize}
        \item 外部融资依赖可抵押资产。
        \item 实物资本兼具生产功能与抵押功能。
        \item 研发更难估值、清算与抵押。
      \end{itemize}
    \column{0.48\textwidth}
      \textbf{数字化冲击}
      \begin{itemize}
        \item 数据留痕改善信息可验证性。
        \item 经营现金流更容易契约化。
        \item 融资契约减少对实物抵押的依赖。
      \end{itemize}
  \end{columns}
  \vspace{0.7em}
  \begin{center}
    \large 预测：抵押囤积型企业收缩 CapEX，研发投入无条件扩张
  \end{center}
\end{frame}

\begin{frame}{理论预测：门槛命题}
  \begin{itemize}
    \item 企业同时面临资产抵押与现金流抵押双重摩擦。
    \item 数字化缓解融资约束后，有形资本调整方向取决于企业抵押能力是否超过临界水平。
    \item 抵押能力低于阈值：原本融资不足，约束缓解后增加有形资本投入。
    \item 抵押能力高于阈值：原本为融资囤积实物资产，约束缓解后收缩冗余产能。
    \item 研发投入扩张不依赖于有形资本调整方向，是本文实证预测的核心不对称性。
  \end{itemize}
\end{frame}
